\documentclass{mipt-thesis-bs}
% Следующие две строки нужны только для biblatex. Для inline-библиографии их следует убрать.
% \usepackage{mipt-thesis-biblatex}
% \addbibresource{example.bib}

\title{Изучение мохнатых существ с $\alpha$ Центавра}
\author{Сидоров А.\,Б.}
\supervisor{Иванов В.\,Г.}
%\referee{Петров Д.\,Е.}       % требуется только для mipt-thesis-ms
\groupnum{082}
\faculty{Факультет радиотехники и компьютерных технологий}
\department{Кафедра }

\begin{document}

\frontmatter
\titlecontents

\mainmatter


\chapter{Введение}

В настоящее время создается множесто язуков программирования, каждый из которых своей целью ставит 
удобство решения определенного класса задач. Большинство ныне существующих языков программирования
относятстя к платформонезависимыми языкам. Главным их преимуществом является высокая переносимость
кода, существует идиома описывающая данный вид поведения(WORA — Write Once, Run Anywhere). Примерами
таких языков являются: ArkTS, C\# Dart, Java. 

Высокая переносимость обеспечивается по средствам специальных сред исполнения, называемых 
виртуальными машинами, которые исполняют байт-код. Байт-код -- промежуточный, платформенно-независимый 
код низкого уровня, получаемый при компиляции исходного текста программы. Он состоит из компактных 
инструкций, которые исполняются не напрямую процессором, а виртуальной машиной 
(например, JVM для Java), что позволяет запускать приложения на Windows, Linux, macOS и других 
платформах без перекомпиляции.

В связи со сказанным выше на первый план выходит проблемя отслеживания корректности(здесь и далее под
корректностью будем понимать соответсвие архитектуре набора команд и семантике языка), 
сгенерированного байт-кода. Программы отслеживающий корректность назывются верификаторами байт-кода.
Важность изучения и разработки данных программы обуславливается тем, что исполняемые приложения 
могут отвечать за функционирование
крайне чувствительных областей, таких как телекоммуникации, банковская сфера, медицинские услуги, 
где ошибка исполнения может привести к необратимому ущербу. Наличие 
верификатора позволяет обнаруживать ошибки до времени исполнения. Кроме того, верификатор 
оказывается полезен на этапе разработки языка программирования.

Существует несколько источников некорректного байт-кода:
\begin{enumerate}
    \item Ошибка компилятора на этапе кодогенерации, Lowering(понижение), регистровой аллокации и т.д.
    \item Повреждения пакета с исполняемым кодом во время его передачи по сети.
\end{enumerate}

------

В данной работе фокус идет на повышение точности работы верификатора байт-кода путем расширения его 
типовой система и увеличения множества верифицируемых инструкций.

\chapter{Постановка проблемы}
\section{Цели работы:}

\begin{enumerate}
    \item Расширить типовую систему верификатора для более точного отображения типовой
    системы виртуальной машины на типовую систему верификатора.
    \item Разработать и внедрить алгоритмы проверок ряда инструкций, для раширения множества
    верифицируемых программ.
\end{enumerate}

\section{Задачи работы:}
\begin{enumerate}
    \item Изучить существующие решения.
    \item Сравнить типовые системы виртуально машины и верификатора, выявить неверифицируемые типы.
    \item Расширить типовую системы и поддержать логику обработки данных типов.
    \item Проанализировать архитектуре набора команд платформы и выявить множество неверифицируемых
    инструкций.
    \item Реализовать и внедрить алгоритм верификации новых инструкций.
    \item Протестировать каждый из предложенных алгоритмов.

\end{enumerate}

Цели работы можно считать достигнутыми. Так как по ходу работы было реализовано порядка 10 обработок
инструкций, что позволило выявить ошибки кодогенерации компилятора и улучшить гарантии платформы 
по безопасности. Как итог ahead-of-time(заранее) режим верификатора был добавлен в основную ветку OHOS 

\chapter{Обзор существующих решений}

На данный момент большинство виртуальных машин обладают верификаторами, в качестве самых популярных 
примеров могут послужить JVM - виртуальная машина для исполнения байт-кода таких языков как: java, kotlin;
.Net CLR - виртуальная машина для исполнения байт-кода С\#, F\#. Так же существую верификаторы, которые
не поставляются вместе с платформой, а являются сторонними библотеками, примерами могут послужить 
BCEL Verifier для Java байт-кода и PEVerify, который проверяет на корректность IL-кода(С\#). 
Рассмотрим некоторые из существующих решений и виды выполняемых проверок более подробно.

\section{Верификация байт-кода статически типизированных яп:}
В статически типизированных языках типы переменных, аргументов и возвращаемых значений известны на 
этапе компиляции. Это позволяет компилятору генерировать байткод с встроенными гарантиями типов.
Но байткод может быть сгенерирован не только компилятором, поэтому верификатор остается неотъемлимой частью
вертуальной машины. Перечислим основные выполняемых проверок:

\begin{enumerate}
    \item \textbf{Проверки формата байткода:} 
        
    Целью данных проверок является тестирование на валидность инструкций, корректность их формата и 
    правильность индексов, соответствующих методам, полям и типам.

    \item \textbf{Проверки графа управления} 

    Во время выполнения данного типа проверок, программа анализируется 
    на корректность инструкций условного перехода: целевой адрес не выходит за пределы 
    текущего метода, в середину try-блока или в середину другой инструкции. Проверяется отсутствие циклов 
    с некорректным состоянием регистров. Контролируется, что все пути выполнения заканчиваются 
    инструкциями обозначающими $return$ или $throw$

    \item \textbf{Проверки типовой безопастности регистров:}
    
    Байткод инспектируется на предмет того, что тип каждого регистра является примитивным или 
    ссылочным. Верифицируется, что любой регист используется после своей инициализации.

    \item \textbf{Проверки инструкций вызова:}
    
    Инструкции вызова проходят проверку на корректность объекта ресивера(в случае нестатического метода).
    На количество аргументов, на типы данных аргументов и тип возвращаемого значения. Верифицируется
    возможность доступа.

    \item \textbf{Проверка доступа обращения:}
    
    Именно данная верификация является одной из самых важных для безопасности. Как известно, 
    существует 3 основных модификатора доступа: $публичные$, $защищенные$ и $приватные$, которые
    применимы, как к элементам класса, так и классам и целым модулям. Контроль того, что их семантика
    не нарушается на уровне байткода необходим.
\end{enumerate}

Рассмотрим пару примеров промышленных решений
\begin{enumerate}

    \item \textbf{ART}
    
    В контексте данной работы для нас наибольший интерес представляют верификаторы тех платформ, которые 
    предназначены для исполнения кода на мобильных устройствах. Именно таким является верификатор ART(Android
    runtime). ART исполняет DEX-байткод, источником которого могут быть: сторонние приложения, 
    сгенерированного кода, потенциально повреждённых или злонамеренных источников. 
    \item \textbf{LLVM} 
    
    LLVM-верификация — это обладает несколько иной философия, чем ART: она не про безопасность 
    пользователя, а про корректность как ПП(здесь и далее ПП - промежуточное представление) компилятора. ПП 
    генерируется доверенным фронтендом, но оптимизации могут его сломать. Можно выделить 4 уровня верификации:
    уровень модудля, функций, базовых блоков и инструкций. Проверяются как локальные, так глобальные инварианты.
\end{enumerate}

\section{Верификация байт-кода динамически типизированных яп:}
В динамически типизированных языках типы определяются в процессе исполнения: переменная может менять тип . 
Байткод здесь часто проще, и верификатор фокусируется не на статических типах, а на структурной целостности 
и совместимости.

% \begin{lstlisting}[language=Python]
% def f(x):
%     return x*x
% \end{lstlisting}

\backmatter

% Здесь идет  текст. Вот так выглядит ссылка на библиографию \cite{langmuir26}. Аналогично добавляются еще главы, внутри них можно объявлять секции с помощью \verb|\section|.
% \printbib
% Следующие строки необходимо раскомментировать, а предыдущую закомментировать, если используется inline-библиография.
%\begin{thebibliography}{99}
%    \bibitem{langmuir26}
%        H. Mott-Smith, I. Langmuir. ``The theory of collectors in gaseous discharges''. \emph{Phys. Rev.} \textbf{28} (1926)
%\end{thebibliography}

\chapter{Благодарности}

Благодарности идут тут.

\end{document}